\section{Related Works}
Financial time-series forecasting has long used recurrent neural networks to model nonlinear temporal dependence in stock and index data. LSTM-type models remain practical baselines because they can encode sequential price behavior without requiring very large architectures. Recent studies also extend this family with bidirectionality, convolutional preprocessing, attention, or evolutionary feature construction, including CNN-BiLSTM-Attention models for stock-index prediction \cite{yang2024cnn_bilstm_attn}, Indonesian blue-chip stock comparisons with LSTM baselines \cite{saleh2025indonesian_lstm}, and LSTM variants combined with symbolic genetic programming for large-scale stock-return forecasting \cite{li2024sgp_lstm_stock}. These works motivate using LSTM/BiLSTM as the primary recurrent benchmark while testing whether recurrent variants add value.

Transformer-based forecasting has become a second major direction for long and multivariate time series. General forecasting architectures such as Temporal Fusion Transformer (TFT) \cite{lim2019tft}, Informer \cite{zhou2020informer}, Autoformer \cite{wu2021autoformer}, and TimeXer \cite{wang2024timexer} introduce attention, sequence distilling, decomposition, or exogenous-variable handling for temporal prediction. Finance-specific work has also adapted self-attention to stock selection and market forecasting, including Stockformer for price-volume factor modeling \cite{ma2024stockformer} and time-series Transformer models with SHAP/LIME analysis on BIST100 banking stocks \cite{calik2025xai_transformer_bist}. These studies provide the basis for treating Informer and Autoformer as supplementary Transformer baselines rather than as replacements for recurrent statistical validation.

External variables are frequently used to enrich financial forecasting when their timing and economic relationship are informative. Prior work has studied macroeconomic indicators for Indonesian stock-market prediction \cite{adyatma2022idx_macro}, macro and sentiment features with LSTM-Attention \cite{begum2024macro_sentiment_lstm}, external indicators for bank-stock forecasting \cite{dechjijaruwate2025thai_banks_external}, and exogenous-variable Transformer designs \cite{qiu2025dag_exogenous,wang2024timexer}. This literature supports an explicit comparison between \texttt{Without\_Macro} and \texttt{With\_Macro} regimes, while leaving open whether macro variables help uniformly across banks.

Interpretability methods are useful for checking which inputs a forecasting model relies on across assets, seeds, and feature regimes. SHAP provides a model-agnostic additive attribution framework \cite{lundberg2017shap} and has been applied to stock-market factor analysis and neural forecasting workflows \cite{shih2024lightgbm_shap,begum2025shap_rnn_stock,gupta2025lstm_shap_xai}. Other recent financial-forecasting studies use LRP, dashboards, or broader explainable-deep-learning workflows \cite{tanveer2024xai_stock,tanveer2024lrp_stock,jagannathan2025hybrid_xai_stock,ayyappan2025xdl_stock}. In this study, SHAP is used as a feature-attribution layer for comparing model reliance patterns, not as a causal explanation of market movement.

Beyond architecture and features, applied financial forecasting often reports limited robustness checks. Single-run comparisons can obscure seed sensitivity, bank-specific behavior, and differences between average error improvements and paired statistical evidence. This paper addresses that gap by combining controlled ablation, repeated-seed evaluation, and paired statistical tests for the original recurrent benchmark, then reporting Transformer and recurrent-variant extensions with explicit caveats about the remaining cross-family validation.
